\labsessionheader{2026-09-24}{Rail-bank follow-ups in four parallel tracks (xdu2.1, .8, .9, .10)}%
  {ddde474f}{n/a}{docs/prd/phase-04-swap-route-control/inc1-records/session-2026-09-24-railbank-parallel/}

\subsection*{Goal}
Clear the four rail-bank follow-ups filed by \bd{test\_C\_libtooling-xdu2.4}.
One orchestrator ran four track agents, each in its own git worktree with its
own tests. Each track then got independent Rule-13 reviewers, and the shared
suites ran only at the orchestrator's merge gate.

\subsection*{What landed}
\begin{itemize}
  \item \commit{d0720cb1} --- Step 0: the rail-bank C matrix reads its
    expectations from a JSON sidecar per source, so each track owns its own
    matrix (null control: same 27 verdicts; 5 of 5 sidecar mutants red).
  \item \commit{33f8bcdb}, \commit{ddde474f} --- \bd{test\_C\_libtooling-xdu2.1}
    (Track B), CLOSED. See Findings.
  \item \commit{bffd2150} --- \bd{test\_C\_libtooling-xdu2.9} (Track D), CLOSED:
    RailBankPromote PASS H resolves a constant emptied-arm condition to its
    taken edge. A constant-dead cell store is no longer a bank; an always-taken
    arm becomes an unconditional bank store.
  \item \commit{2ce11e36} --- \bd{test\_C\_libtooling-xdu2.8} (Track C), stays
    OPEN: PredMux E4 arm (1d) admits a classical runtime-VLA read whose bounds
    proof is local to its arm. The load is rewritten to
    \texttt{select(idx <u extent, \&p[idx], fresh zero cell)}. The bead's own
    program still declines at a second, pre-existing wall.
  \item \commit{ea5eaee3} --- \bd{test\_C\_libtooling-xdu2.10} (Track A), CLOSED:
    the carve's C6 read-order test asks whether the read executing implies the
    write loop ran, using classical facts and \texttt{isImpliedCondition}.
    Before, it asked only for dominance.
  \item \commit{b3f51e0a} --- the matrix runner refuses duplicate sidecar keys.
\end{itemize}

\subsection*{Findings}
\begin{itemize}
  \item \fconfirmed\ \fpreexisting\ \bd{xdu2.1} was already closed for the
    \texttt{cq\_theta} relabel by \commit{95949550}, but reviewers found two
    more P1s of the same class. (r2) A rail parked \emph{before} an in-loop
    in-place \texttt{xor}/\texttt{add} was read after it: the fence took the
    park store as the read. (r3) A rail parked \emph{after} the in-loop op made
    every cell read the last trip's value, because the fence cut at the SSA
    ``refresh'', which is wrong for a memory cell. Both are fenced (TRACK\_B.md).
  \item \fconfirmed\ Track A, round 2 (P2, latent): tainted branch conditions
    were used as facts. Round 3 (P2): the fix to that over-declined the Grover
    shape. Both fixed; round 4 clear.
  \item \fconfirmed\ Track C, round 2 (P2): \texttt{cq-opt} aborted on a
    non-zero address space; round 3 (P3): it reused a forged zero-cell global.
    Both fixed, the second by never reusing a cell.
  \item \foverdecline\ Filed: \bd{xdu2.11}, \bd{xdu2.12} (dead-cell spellings),
    \bd{xdu2.14} (in-loop guarded-read wall), \bd{xdu2.15} (park-after-op
    twin), \bd{yiid.1} (bank/register twin disagreement).
\end{itemize}

\subsection*{Measurements}
\begin{tabular}{@{}lrl@{}}\toprule
Quantity & Value & Command \\ \midrule
ctest entries (all pass, final merge) & 491 & \measured{ctest --test-dir build -j6} \\
lit fixtures & 2232 & \measured{ls ir-pass/test/*.ll | wc -l} \\
cq-opt equivalence mismatches & 0 of 2232 & \measured{--target cq-opt-equivalence-full} \\
corpus A/B, B merge: A$\to$D / D$\to$A / stream & 18 / 0 / 0 & \measured{session ab.py, 2802 modules, null 0} \\
corpus A/B, A merge: A$\to$D / D$\to$A / stream & 0 / 15 / 0 & \measured{session ab.py, 2772 modules, null 0} \\
corpus A/B, C merge: A$\to$D / D$\to$A / stream & 0 / 18 / 0 & \measured{session ab.py, 2739 modules, null 0} \\
corpus A/B, D merge: A$\to$D / D$\to$A / stream & 0 / 14 / 0 & \measured{session ab.py, 2692 modules} \\
Rule-13 rounds (A / B / C / D) & 4 / 4 / 3 / 2 & independent reviewers \\
\bottomrule\end{tabular}

In each A/B, every mover is one of that track's own new fixtures or matrix rows.
The 18 A$\to$D at the B merge are the P1 closures.

\subsection*{Verification run}
Every merge: rebuild, the track's own tests, full \texttt{ctest -j6},
\texttt{cq-opt-equivalence-full}, corpus A/B against the pre-merge binary (lit
fixtures, stage 1 of every e2e \texttt{.c}, every matrix case, and
\fpath{example/grover/main.c}), the exact Grover run, and the CODE\_MAP check.
Grover stayed at exit 0 with 0 \texttt{cqrt\_qram\_*} and STRICT 2b 3/0/0
throughout. The D-merge A/B's null control ran after the A/B, not first. At
Step 0, full ctest was 485/486: \texttt{code\_map\_drift} failed on drift
already present at \commit{2143a535}, and the regenerated map fixed it. Not
run: \texttt{+Asserts}, multi-function-module probes (\bd{1yeu}).

\subsection*{Open questions and next steps}
\begin{itemize}
  \item Owner decision \bd{xdu2.13}: how to lower a 4.4 data select whose
    condition mixes a classical and a quantum conjunct. It blocks
    \bd{xdu2.8}'s program.
  \item Wave 2 is not started: \bd{xdu2.7}, \bd{xdu2.5}, \bd{xdu2.6}.
\end{itemize}

\subsection*{Errata to earlier entries}
None.
